% ==== P4 draft O1 — §15.1 대체 블록 ====
\subsection{왜 흑연에는 없고 LCO 에는 있는가 --- MIT 의 물리}\label{sec:lco-why}
삽입 반응의 전체 엔트로피 변화는 ``리튬 자리 배열''(config)$\cdot$``격자 진동''(vib)$\cdot$``전자 준위 점유''(electronic)
세 자유도의 합이며, 흑연에서는 셋째 몫이 작다 --- 흑연$\cdot$$\mathrm{LiC_6}$ 모두 전도성이 좋아 충방전 동안 $g(E_F)$ 가
크게 바뀌지 않아 전자 분포의 엔트로피 \emph{변화}가 거의 없다. LCO 는 다르다.
$x{=}1$($\mathrm{LiCoO_2}$)에서 Co$^{3+}$ 는 팔면체 결정장에서 $t_{2g}^6$ low-spin 닫힌 껍질을 이루고, 이 가득 찬
$t_{2g}$ 띠가 그 위 빈 띠와 갭으로 갈려 Fermi 준위가 갭 안에 놓인다 --- 전자간 상호작용 없이 띠 채움만으로 절연이
서는 \emph{밴드 절연체}(band insulator)이고, 이것이 전기적 절연($g(E_F)\!\approx\!0$)의 물리 부류다\cite{menetrier1999,motohashi2009}.
탈리튬화로 $x$ 가 $\sim$0.94 아래로 내려가면 Co$^{4+}$ 정공이 생겨 $t_{2g}$ 띠에 전도 전자가 열려 \emph{금속}이
된다($g(E_F)\!\to\!$유한).

다만 절연$\to$금속의 문턱은 단순한 밴드 채움 논리로는 닫히지 않는다. 정공이 부분적으로 빈 띠를 만들면 곧바로
금속이어야 하나, 실측 전이는 $x\approx0.75$--$0.94$ 의 2상 공존을 동반한 \emph{1차 전이}이며, 정공이 즉시
비국재화하지 않고 임계 조성까지 국재해 있다가 한꺼번에 열린다\cite{reimers1992,motohashi2009} --- 이 거동은
단일입자 밴드 그림 밖의 전자 상관(correlation) 물리를 요구한다. 그럼에도 전자 분포가 새로 켜지는 곳은 이 전이의
2상역(\S\ref{sec:lco-hys} 의 T1)에 국한되므로, 전자 준위 점유의 엔트로피 \emph{변화}도 그 구간에 국소한다.

\begin{bgbox}
\textbf{금속--절연체 전이와 Mott 물리 ---} 절연체는 미시 기원에 따라 두 부류로 갈린다. 하나는
\emph{밴드 절연체}(band insulator)로, 단일입자 띠 그림에서 원자가띠가 가득 차고 그 위 전도띠와 갭으로 갈려
Fermi 준위가 갭 안에 놓이는 경우이며, 전자간 상호작용을 넣지 않은 띠 채움만으로 절연이 설명된다. 다른 하나는
\emph{Mott 절연체}(Mott insulator)로, 띠 그림으로는 띠가 반만 차 금속이어야 하나 전자간 Coulomb 반발 $U$ 가
자리 사이 이동적분 $t$ 를 눌러 전자가 제자리에 국재하는 경우이며, 띠 채움이 아니라 전자 상관(correlation)이
절연을 만든다\cite{mott1968,imada1998}. 금속--절연체 전이(MIT)는 채움(도핑)$\cdot$띠폭(압력)$\cdot$상관의
경쟁으로 일어나고, 어느 축이 지배하느냐가 전이의 성격을 정한다\cite{imada1998}. $\mathrm{LiCoO_2}$($x{=}1$)에
이를 적용하면, $t_{2g}^6$ 닫힌 껍질의 가득 찬 $t_{2g}$ 띠는 \emph{밴드 절연체}이지 Mott 절연체가 아니다\cite{menetrier1999}.
탈리튬화가 넣는 정공은 밴드 그림이면 부분적으로 빈 띠로서 즉각적인 금속화를 예상시키나, 실측 MIT 는
$x\approx0.75$--$0.94$ 2상역의 \emph{1차 전이}로 나타난다\cite{menetrier1999,reimers1992,motohashi2009}.
Marianetti, Kotliar, Ceder 는 이 1차 전이를, 정공이 Li 빈자리에 속박되어 갭 안 \emph{불순물 준위}를 이루고 그
불순물 띠가 임계 정공 농도에서 비국재화하는 \emph{불순물 밴드의 Mott 전이}로 설명했다\cite{marianetti2004}.
상관 물리가 host $t_{2g}$ 띠가 아니라 불순물 준위에서 작동한다는 것이 이 그림의 핵심이며, host 띠 자체를
Mott 절연체로 보는 주장은 아니다. forward 모델은 이 미시 기작 전체를 조성축 게이트 $g(E_F,x)$ 하나로
현상학화한다 --- 게이트 중심 $x_\mathrm{MIT}$ 은 실측 2상역 중앙의 전이 조성으로, 폭 $\Delta x_\mathrm{MIT}$ 은
결함이 흩뜨리는 발현 조성의 분산으로 대응하며, 그 게이트 형태의 정당화는 \S\ref{sec:lco-gate} 가 닫는다. 미시
기작이 \emph{왜} 절연이 임계 조성까지 버티는가를 답한다면, 현상학 게이트는 그 순 효과인 $g(E_F)$ 의 켜짐을
\emph{어떻게} 하나의 매끄러운 조성 함수로 담는가를 답한다.
\end{bgbox}

이렇게 전자항이 켜지는 게이트는 MIT 의 직접 결과이며, 다음 소절 \S\ref{sec:lco-Se} 가 이 켜짐을
Fermi--Dirac 분포에서 Sommerfeld 전자 엔트로피로 유도해 전이당(삽입당) 전자 엔트로피 몫을 닫는다.
% ==== 블록 끝 ====
%
% NOTE — Read Coverage / 보존 확인 / 자체검수
% [Read Coverage] AUTHOR_BRIEF.md 전문; ch1_sec15_lcoelec.tex 전문(§15.1 원문 L12–20 + §15.4 sec:lco-gate
%   게이트-형태 정당화 L146–193 — 역할분담: 본 bgbox=왜(기작), §15.4=게이트 형태 정당화, 중복 회피);
%   ch1_sec13_lcohys.tex L120–160(T1 sec:lco-hys-mit 슬롯 분리 — 전자 엔트로피는 ∂U/∂T 슬롯, 히스 Ω gap 과
%   별개이고, 미시 기작·전자 DOF 배경을 sec:lco-why 로 위임한다는 지시 확인); V1020_STYLE_RUBRIC.md 전문;
%   ch1_sec02a_part0.tex L171–188(house bgbox 양식); ch1_sec11 L42(MIT 병기 최초 정의 확인).
% [보존 확인] (1) config·vib·electronic 세 자유도 합·흑연 셋째 몫 작음(g(E_F) 미변) 문장 전량 보존;
%   (2) x=1 Co3+ t2g^6 low-spin 닫힌 껍질 절연 + \cite{menetrier1999,motohashi2009} 보존(밴드 절연체 부류로 정밀화);
%   (3) 탈리튬화 x<~0.94 Co4+ 정공→금속(g→유한) 문장 보존; (4) MIT 2상역 T1(\S\ref{sec:lco-hys}) 국소화 문장 보존;
%   (5) "전자항이 켜지는 게이트 = MIT 의 직접 결과" 취지 보존(마무리 다리); (6) §15.2 유도(sec:lco-Se) 로의 다리 유지.
% [자체검수 — 물리 하드가드 3항] (G1) x=1 LiCoO2 = 밴드 절연체로만 서술, "Mott 절연체" 지칭 0
%   (+ "host t2g 띠를 Mott 절연체로 보는 주장 아님" 오독 방지 가드; A2/D2 대조형 예외 적용); (G2) Marianetti–Kotliar–Ceder
%   = "불순물 밴드의 Mott 전이로 설명했다" 수준 서술, 상관 물리를 불순물 준위 소관으로 국한, "확정된 유일 기작" 단정 회피;
%   (G3) menetrier1999·motohashi2009 유실 없음; 신규 수식·라벨 0(인라인 기호만, 기존 label 재사용); 인용 6키(mott1968,
%   imada1998, marianetti2004, menetrier1999, reimers1992, motohashi2009)만 사용. §13 T1 슬롯 분리와 무모순
%   (전자 엔트로피를 구조적 히스 Ω gap 에 섞지 않고, ∂U/∂T 슬롯의 게이트 g(E_F,x) 로만 전달).
